\section{Sinh Quỹ đạo Kinodynamic}\label{sec:kinodynamic_theory}

Lập kế hoạch chuyển động hình học (geometric motion planning) tìm kiếm một đường đi
không va chạm trong không gian cấu hình của robot mà không xét đến các đặc tính vật lý
của hệ thống.
Tuy nhiên, trong thực tế, robot là một hệ cơ học mà trạng thái của nó chỉ có thể tiến triển
theo thời gian theo các phương trình vi phân xác định được dẫn dắt bởi các đầu vào cơ cấu
chấp hành.
Một đường đi không va chạm về mặt hình học vẫn có thể không khả thi về mặt vật lý nếu nó
đòi hỏi, chẳng hạn, sự thay đổi tức thời của vận tốc hoặc các lực vượt quá giới hạn bão hòa
của cơ cấu chấp hành.
Lập kế hoạch chuyển động kinodynamic mở rộng bài toán hình học bằng cách yêu cầu quỹ đạo
được lập kế hoạch phải đồng thời tránh vật cản, thỏa mãn các ràng buộc vi phân do động lực
học hệ thống áp đặt, và tối ưu hóa một tiêu chí hiệu năng do người dùng chỉ định.

Mục này xây dựng bài toán lập kế hoạch kinodynamic dưới dạng toán học chặt chẽ
(Mục~\ref{subsec:kinodynamic_problem}), giới thiệu tham số hóa quỹ đạo đa thức bằng
đường cong Bézier như một kỹ thuật để đưa bài toán về dạng hữu hạn chiều
(Mục~\ref{subsec:bezier_parameterization}), và trình bày các ràng buộc liên tục đảm bảo
chuyển động trơn qua các điểm chuyển tiếp giữa các đoạn quỹ đạo
(Mục~\ref{subsec:continuity_constraints}).

\subsection{Bài toán Lập kế hoạch Kinodynamic}\label{subsec:kinodynamic_problem}

Xét một hệ thống robot có trạng thái $x(t) \in \mathcal{X} \subseteq \mathbb{R}^{n_x}$
và đầu vào điều khiển $u(t) \in \mathcal{U} \subseteq \mathbb{R}^{n_u}$ tiến triển theo
phương trình vi phân thường
\begin{equation}
    \dot{x}(t) = f\!\left(x(t),\, u(t)\right), \quad t \in [0, T],
    \label{eq:kinodynamic_dynamics}
\end{equation}
trong đó $f \colon \mathbb{R}^{n_x} \times \mathbb{R}^{n_u} \to \mathbb{R}^{n_x}$ là ánh xạ
động lực học của hệ và $T > 0$ là độ dài chân trời lập kế hoạch, có thể được cố định trước
hoặc được coi là biến tối ưu.
Vector trạng thái thường mã hóa các vị trí và vận tốc suy rộng của robot, và tùy theo
mô hình động lực học có thể bao gồm thêm các đạo hàm bậc cao hơn.

Một quỹ đạo kinodynamic là một cặp $\left(x(\cdot), u(\cdot)\right)$ xác định trên $[0,T]$
thỏa mãn phương trình động lực học~\eqref{eq:kinodynamic_dynamics} cùng với các yêu cầu sau.

\begin{enumerate}
    \item \textbf{Điều kiện biên.}
          Quỹ đạo phải xuất phát từ trạng thái đầu được chỉ định và đến được vùng đích mong muốn:
          \begin{equation}
              x(0) = x_{\mathrm{start}}, \qquad x(T) \in \mathcal{X}_{\mathrm{goal}},
              \label{eq:boundary_conditions}
          \end{equation}
          trong đó $\mathcal{X}_{\mathrm{goal}} \subseteq \mathcal{X}$ là tập đích.

    \item \textbf{Ràng buộc an toàn.}
          Quỹ đạo trạng thái phải nằm trong vùng không có vật cản
          $\mathcal{X}_{\mathrm{free}} \subseteq \mathcal{X}$ tại mọi thời điểm:
          \begin{equation}
              x(t) \in \mathcal{X}_{\mathrm{free}}, \quad \forall\, t \in [0, T].
              \label{eq:safety_constraint}
          \end{equation}

    \item \textbf{Giới hạn cơ cấu chấp hành và động học.}
          Đầu vào điều khiển phải nằm trong tập khả thi của cơ cấu chấp hành:
          \begin{equation}
              u(t) \in \mathcal{U}, \quad \forall\, t \in [0, T].
              \label{eq:control_constraint}
          \end{equation}
          Trong nhiều bài toán thực tế, $\mathcal{U}$ được đặc trưng bởi các giới hạn tường minh
          lên vận tốc, gia tốc và độ giật dọc theo quỹ đạo,
          ví dụ\ $\|\dot{x}(t)\| \le v_{\max}$, $\|\ddot{x}(t)\| \le a_{\max}$
\end{enumerate}

Mục tiêu là tìm một quỹ đạo tối thiểu hóa phiếm hàm chi phí có dạng
\begin{equation}
    J\!\left(x(\cdot), u(\cdot)\right)
    = \int_{0}^{T} L\!\left(x(t), u(t)\right) \mathrm{d}t + E\!\left(x(T)\right),
    \label{eq:kinodynamic_cost}
\end{equation}
trong đó $L$ là running cost phạt, chẳng hạn, nỗ lực điều khiển hoặc độ lệch so với
quỹ đạo tham chiếu, và $E$ là terminal cost.

Có hai khó khăn cấu trúc phân biệt bài toán này với bài toán hình học thuần túy.
Thứ nhất, ràng buộc an toàn~\eqref{eq:safety_constraint} phải được thỏa mãn liên tục
trên toàn bộ khoảng $[0, T]$, không chỉ tại một tập hữu hạn các điểm rời rạc,
và $\mathcal{X}_{\mathrm{free}}$ nói chung là phi lồi do sự hiện diện của vật cản.
Thứ hai, phiếm hàm chi phí~\eqref{eq:kinodynamic_cost} và ràng buộc động lực
học~\eqref{eq:kinodynamic_dynamics} có tính chất vô hạn chiều, vì các biến quyết định
$x(\cdot)$ và $u(\cdot)$ là các hàm theo thời gian.
Để giải bài toán này bằng phương pháp số, cần phải đưa nó về dạng bài toán hữu hạn chiều
bằng cách chọn một tham số hóa quỹ đạo phù hợp.

\subsection{Tham số hóa Quỹ đạo bằng Đường cong Bézier}\label{subsec:bezier_parameterization}

Một hướng tiếp cận tự nhiên để tham số hóa hữu hạn chiều là biểu diễn quỹ đạo dưới dạng
đa thức với các hệ số đóng vai trò là biến quyết định.
Trong số các biểu diễn đa thức, đường cong Bézier đặc biệt phù hợp với lập kế hoạch
kinodynamic vì các tính chất hình học và giải tích của chúng cho phép áp đặt các ràng buộc
an toàn, độ trơn và tính khả thi động học trực tiếp dưới dạng các điều kiện đại số trên
các hệ số.

\subsubsection{Định nghĩa}

Một đường cong Bézier bậc $d$ trên khoảng tham số $[0,1]$ được xác định bởi
$d+1$ \emph{điểm điều khiển} $\gamma_0, \ldots, \gamma_d \in \mathbb{R}^{n}$
thông qua các đa thức cơ sở Bernstein.
Đa thức Bernstein thứ $k$ bậc $d$ được định nghĩa là
\begin{equation}
    \beta_{k,d}(s) := \binom{d}{k} s^{k}(1-s)^{d-k},
    \quad k = 0, \ldots, d, \quad s \in [0,1].
    \label{eq:bernstein}
\end{equation}
Theo định lý nhị thức, các đa thức này không âm và phân hoạch đơn vị:
$\sum_{k=0}^{d} \beta_{k,d}(s) = 1$ với mọi $s \in [0,1]$.
Do đó, với mỗi $s$ cố định, các giá trị $\{\beta_{k,d}(s)\}_{k=0}^{d}$ tạo thành một
tập hệ số tổ hợp lồi.
Đường cong Bézier được định nghĩa là trung bình có trọng số tương ứng của các điểm điều khiển:
\begin{equation}
    \gamma(s) := \sum_{k=0}^{d} \beta_{k,d}(s)\,\gamma_k,
    \quad s \in [0,1].
    \label{eq:bezier_curve}
\end{equation}

\subsubsection{Các tính chất liên quan đến lập kế hoạch quỹ đạo}

Ba tính chất của đường cong Bézier được khai thác trực tiếp trong lập kế hoạch quỹ đạo
kinodynamic.

\textbf{Nội suy điểm đầu cuối.}
Đường cong đi qua chính xác điểm điều khiển đầu tiên và điểm điều khiển cuối cùng:
\begin{equation}
    \gamma(0) = \gamma_0, \qquad \gamma(1) = \gamma_d.
    \label{eq:endpoint_interpolation}
\end{equation}
Tính chất này cho phép áp đặt các điều kiện biên về vị trí và, thông qua công thức đạo hàm
bên dưới, cả về vận tốc, dưới dạng các ràng buộc đẳng thức tuyến tính trên các điểm điều khiển.

\textbf{Tính chất bao lồi.}
Toàn bộ đường cong nằm bên trong bao lồi của các điểm điều khiển của nó:
\begin{equation}
    \gamma(s) \in \operatorname{conv}\!\left(\{\gamma_k\}_{k=0}^{d}\right),
    \quad \forall\, s \in [0,1].
    \label{eq:convex_hull}
\end{equation}
Suy ra rằng nếu tất cả $d+1$ điểm điều khiển nằm trong một tập lồi $\mathcal{S}$,
thì $\gamma(s) \in \mathcal{S}$ với mọi $s$.
Nguyên tắc bao hàm này chuyển đổi ràng buộc an toàn liên tục~\eqref{eq:safety_constraint}
thành một tập hữu hạn các ràng buộc thuộc tập trên các điểm điều khiển, với điều kiện vùng
an toàn có thể được xấp xỉ hoặc phân hoạch thành các tập lồi.

\textbf{Công thức đạo hàm (Hodograph).}
Đạo hàm của $\gamma$ theo tham số $s$ cũng là một đường cong Bézier bậc $d-1$,
với các điểm điều khiển được tạo thành từ các sai phân hữu hạn bậc một của các điểm
điều khiển ban đầu:
\begin{equation}
    \dot{\gamma}(s) = \sum_{k=0}^{d-1} \beta_{k,d-1}(s)\,\dot{\gamma}_k,
    \qquad \dot{\gamma}_k = d\!\left(\gamma_{k+1} - \gamma_k\right).
    \label{eq:bezier_derivative}
\end{equation}
Áp dụng đệ quy~\eqref{eq:bezier_derivative} cho phép tính gia tốc và các đạo hàm bậc cao
hơn dưới dạng đường cong Bézier có bậc giảm dần, với các điểm điều khiển được biểu diễn
dưới dạng tổ hợp tuyến tính của các $\gamma_k$ ban đầu.
Do đó, các ràng buộc lên vận tốc và gia tốc dọc theo quỹ đạo --- chẳng hạn các giới hạn
động học trong~\eqref{eq:control_constraint} --- quy về các ràng buộc bất đẳng thức tuyến
tính trên các sai phân điểm điều khiển, và có thể được áp đặt đồng thời với tính chất bao lồi.

\subsection{Ràng buộc Độ trơn tại Điểm chuyển tiếp}\label{subsec:continuity_constraints}

Với các quỹ đạo trải dài qua nhiều vùng không gian làm việc có tính chất động học khác nhau,
một đoạn Bézier duy nhất bậc vừa phải có thể không đủ linh hoạt.
Biện pháp thực tế là sử dụng quỹ đạo \emph{Bézier từng đoạn} (piecewise Bézier) gồm $M$
đoạn ghép nối đầu-đuôi.
Gọi đoạn thứ $i$ là đường cong Bézier bậc $d$ được xác định bởi các điểm điều khiển
$a^{(i)}_0, \ldots, a^{(i)}_d$, với $i = 1, \ldots, M$.
Việc ghép nối các đoạn đảm bảo đường đi liên thông, nhưng cần thêm các điều kiện đại số
tại mỗi điểm nối để đảm bảo tính khả vi của quỹ đạo tổng thể.

\textbf{Liên tục $C^0$ (vị trí).}
Điểm cuối của đoạn $i$ trùng với điểm đầu của đoạn $i+1$:
\begin{equation}
    a^{(i)}_d = a^{(i+1)}_0.
    \label{eq:c0_continuity}
\end{equation}

\textbf{Liên tục $C^1$ (vận tốc).}
Vận tốc ra của đoạn $i$ khớp với vận tốc vào của đoạn $i+1$,
điều này theo~\eqref{eq:bezier_derivative} được dịch thành đẳng thức của các sai phân bậc
một tại biên:
\begin{equation}
    a^{(i)}_d - a^{(i)}_{d-1} = a^{(i+1)}_1 - a^{(i+1)}_0.
    \label{eq:c1_continuity}
\end{equation}

\textbf{Liên tục $C^2$ (gia tốc).}
Gia tốc ra của đoạn $i$ khớp với gia tốc vào của đoạn $i+1$,
được biểu diễn qua các sai phân bậc hai tại biên:
\begin{equation}
    a^{(i)}_d - 2a^{(i)}_{d-1} + a^{(i)}_{d-2}
    = a^{(i+1)}_2 - 2a^{(i+1)}_1 + a^{(i+1)}_0.
    \label{eq:c2_continuity}
\end{equation}

Các điều kiện~\eqref{eq:c0_continuity}--\eqref{eq:c2_continuity} là các ràng buộc đẳng thức
tuyến tính trên các điểm điều khiển và do đó bảo toàn cấu trúc đại số của bài toán tối ưu.
Việc yêu cầu liên tục $C^1$ hay $C^2$ phụ thuộc vào bậc của mô hình động lực học robot:
mô hình tích phân đôi (double-integrator) chỉ cần $C^1$, trong khi mô hình có đầu vào mô-men
xoắn tường minh thường đòi hỏi $C^2$ để tránh các lệnh lực xung tại điểm nối giữa các đoạn.

Sau khi đặt tham số hóa Bézier từng đoạn, bài toán lập kế hoạch kinodynamic quy về bài toán
tối ưu hữu hạn chiều: các biến quyết định là các điểm điều khiển $\{a^{(i)}_k\}$ và, khi áp
dụng, chân trời $T$; hàm mục tiêu là phiếm hàm chi phí~\eqref{eq:kinodynamic_cost} được biểu
diễn qua các điểm điều khiển; và các ràng buộc bao gồm điều kiện
biên~\eqref{eq:boundary_conditions}, các ràng buộc an toàn suy từ tính chất bao
lồi~\eqref{eq:convex_hull}, các giới hạn động học được áp đặt qua công thức đạo
hàm~\eqref{eq:bezier_derivative}, và các điều kiện tại điểm
nối~\eqref{eq:c0_continuity}--\eqref{eq:c2_continuity}.

Cấu trúc này --- hàm mục tiêu tích phân, ràng buộc vi phân trên quỹ đạo, điều kiện biên cố
định, và các ràng buộc bất đẳng thức trên đường --- chính xác là cấu trúc của một Bài toán
Điều khiển Tối ưu (Optimal Control Problem, OCP).
Mục tiếp theo hệ thống hóa mối liên hệ này và phát triển khung OCP làm nền tảng cho phương
pháp tối ưu hóa quỹ đạo được áp dụng trong đồ án này.
